\chapter{测试与评估}
\label{sec:testing}

\section{测试目标与测试范围}

本章围绕 OralSage 的功能完整性、智能体协同、知识检索、资源生成、学习测评、数字问诊、证据一致性和核心响应时间开展测试。测试覆盖学习者从进入系统、建立学习画像、输入学习目标、生成学习计划，到获得课程讲解、进行即时答疑、完成练习并更新学习状态的主要过程，也覆盖数字人患者问诊、自然语言诊断理解、判官接管和临床表达评分。

测试采用自动化测试与真实功能链路测试相结合的方式。自动化测试用于验证后端服务、前端交互逻辑、状态转换和异常处理；真实功能测试使用项目本地知识库、知识图谱、稀疏索引、稠密向量索引和模型服务，验证课程规划、讲义生成、Tutor 讲解和 QA 答疑等模块的实际运行结果。

\begin{table}[H]
  \centering
  \caption{测试环境与测试范围}
  \label{tab:test-environment}
  \begin{tabularx}{\linewidth}{p{0.25\linewidth}X}
    \toprule
    测试项目 & 配置与范围 \\
    \midrule
    运行环境 & Windows 本地运行环境 \\
    后端服务 & OralSage 本地服务，测试地址为 \texttt{127.0.0.1:8000} \\
    自动化测试 & Python 后端测试和 Node.js 前端逻辑测试 \\
    知识检索 & 本地 KAG 检索服务、稀疏倒排索引和 BGE-M3 稠密向量索引 \\
    检索语料 & 34,805 条统一文本片段 \\
    测试主题 & 以“牙周炎”等口腔医学主题为主要测试对象 \\
    测试模块 & 学习画像、学习计划、课程规划、资源生成、Tutor 讲解、QA 答疑、数字问诊、学习测评和状态更新 \\
    \bottomrule
  \end{tabularx}
\end{table}

\section{自动化测试结果}

Python 后端测试共执行 961 项，覆盖学习路由、课程规划、知识检索、学习计划、画像状态、资源生成、即时答疑、测评、缓存和状态持久化等模块，全部完成。

数字问诊功能新增 49 项专项测试，涉及临床训练服务、学习画像事件和实时数字人配置，全部通过。该组测试独立验证患者 Agent、判官 Agent、病例仓库、讯飞语音听写配置、数字人鉴权与画像反馈，便于对新增链路进行集中复核。

前端测试覆盖学习计划草稿加载、计划项目采用、依赖关系处理、主对话状态、学习会话参数和计划入口。其中，学习计划草稿测试完成 60 项，主对话计划草稿测试完成 76 项，主对话测试完成 51 项，学习计划逻辑测试完成全部测试项。

\begin{table}[H]
  \centering
  \caption{自动化测试结果}
  \label{tab:automated-test}
  \begin{tabularx}{\linewidth}{p{0.30\linewidth}p{0.28\linewidth}X}
    \toprule
    测试项目 & 实际结果 & 覆盖内容 \\
    \midrule
    Python 后端测试 & 961 项完成且无失败项 & 路由、规划、检索、资源生成、答疑、测评、画像、缓存与状态管理 \\
    数字问诊专项测试 & 49 项完成且无失败项 & 病例选题、患者交互、诊断识别、判官评分、语音、数字人、状态与画像事件 \\
    学习计划草稿测试 & 60 项完成且无失败项 & 草稿加载、课程范围、推荐卡片、采用请求和异常处理 \\
    主对话计划草稿测试 & 76 项完成且无失败项 & 计划草稿保存、会话切换、计划采用、依赖关系和页面参数 \\
    主对话测试 & 51 项完成且无失败项 & 学习请求、文献请求、答疑、会话状态和学习计划入口 \\
    学习计划逻辑测试 & 测试项全部完成 & 计划生成、计划项目数据和前端状态转换 \\
    \bottomrule
  \end{tabularx}
\end{table}

\section{核心功能测试}

核心功能测试以真实学习流程为对象，验证用户输入是否能够转化为课程规划、学习资源和学习状态。

输入“牙周炎”后，系统生成 5 个课程节点，并为当前课程节点绑定知识节点和资料片段。后续讲义、Tutor 和 QA 均基于该课程节点继续执行。

\begin{table}[H]
  \centering
  \caption{核心功能测试结果}
  \label{tab:function-test}
  \begin{tabularx}{\linewidth}{p{0.27\linewidth}X}
    \toprule
    测试功能 & 实际结果 \\
    \midrule
    学习画像 & 初始画像包含本科高年级、期末考试、病例带入和标准节奏等字段；提问和答题后，薄弱点由空集更新为“机制链条不完整、牙槽骨吸收、牙龈炎与牙周炎区别” \\
    主题学习 & 输入“牙周炎”，生成 5 个课程节点 \\
    学习计划 & 生成计划草稿，保存课程范围、计划项目、日期和依赖关系 \\
    计划反馈修改 & 草稿能够被读取、修改并重新进入计划采用流程 \\
    文献检索 & 输入文献主题后返回对应检索结果 \\
    课程讲解 & Tutor 返回 \texttt{llm\_generated} 内容 \\
    当前章节答疑 & 返回 \texttt{llm\_generated} 内容，未推进课程进度 \\
    课程延伸答疑 & 返回 \texttt{llm\_generated} 内容，并生成 2 条定向检索计划 \\
    学习测评 & 完成题目生成、答案提交、批改、解析和掌握度更新 \\
    学习状态 & 完成课程进度、学习事件、答案记录和掌握度记录保存 \\
    \bottomrule
  \end{tabularx}
\end{table}

\section{智能体协同链路测试}

主题学习流程中的智能体调用顺序为：Planner 根据主题、学习者画像和检索证据生成课程结构；Resource Factory 根据当前课程节点生成讲义；Tutor 根据讲义和证据生成讲解；QA 根据问题范围选择当前证据或发起定向检索；Assessment 负责练习与阶段测评；Memory/User Model 保存学习事件、掌握度和薄弱知识点。

\begin{table}[H]
  \centering
  \caption{智能体协同链路测试结果}
  \label{tab:agent-chain-test}
  \begin{tabularx}{\linewidth}{p{0.25\linewidth}X}
    \toprule
    测试环节 & 实际结果 \\
    \midrule
    Planner & 输入“牙周炎”，生成 5 个课程节点，并为首节点绑定 1 个知识节点和 2 条资料片段 \\
    Resource Factory & 生成首节讲义，返回 177 字概述和 4 个主要知识点 \\
    Tutor & 基于首节讲义生成 \texttt{llm\_generated} 讲解内容 \\
    QA 当前章节模式 & 问题范围为 \texttt{in\_current\_section}，使用 2 条证据直接回答 \\
    QA 延伸检索模式 & 问题范围为 \texttt{course\_related\_but\_not\_current}，生成 2 条检索计划并使用 4 条证据回答 \\
    Assessment & 完成错误答案识别、解析、错因信息和掌握度更新 \\
    Memory/User Model & 答题后更新薄弱点、错误模式和掌握度记录 \\
    Orchestrator & 完成 Planner、Resource Factory、Tutor、QA、Assessment 和 Memory/User Model 的连续调度 \\
    \bottomrule
  \end{tabularx}
\end{table}

两类答疑的 \texttt{advance\_progress} 均为 \texttt{False}，说明用户在课程中插入问题不会导致课程进度被错误推进。

\section{数字问诊专项测试}

专项测试围绕第~\ref{sec:digital-consultation}章的完整流程展开。病例仓库测试检查训练字段完整性、复诊记录排除、直接身份信息清理，以及先选疾病类别再抽取病例的随机选题顺序。患者 Agent 测试使用不同语气的自然语言诊断和追问，检查意图理解、事实编号引用、未知信息处理与疾病名称隐藏，确认患者回复不会直接暴露目标诊断。

患者与判官交接测试覆盖明确诊断识别、诊断复核、用户确认和首道考核题生成。判官测试检查五个考核领域的检索覆盖、题目证据集合、分值范围、批评意见结构和画像写入，非法证据编号无法进入有效评分结果。请求编号测试重复提交会话创建、问诊轮次和考核回答，确认问题序号、分数与学习事件只更新一次。

语音与数字人测试检查讯飞听写签名地址不会泄露 APISecret，医学术语整理会保留学生原始诊断语义，语音输入完成后能够进入患者 Agent。浏览器交互测试还覆盖消息即时显示、输入框自动高度和音频末帧提交，数字人公开配置仅包含建立实时互动所需的参数。49 项专项测试全部通过，结果如表~\ref{tab:clinical-speaking-test} 所示。

\begin{table}[H]
  \centering
  \caption{数字问诊专项测试结果}
  \label{tab:clinical-speaking-test}
  \begin{tabularx}{\linewidth}{p{0.28\linewidth}X}
    \toprule
    检查类别 & 已验证行为 \\
    \midrule
    病例与选题 & 完整字段过滤、脱敏输出、疾病类别随机选择和画像目标映射 \\
    患者交互 & 自然语言意图理解、病例事实约束、未知信息表达和诊断隐藏 \\
    诊断交接 & 明确诊断判断、独立复核、用户确认和判官接管 \\
    判官评分 & 五域证据检索、问题结构、评分范围、批评意见和合法证据引用 \\
    状态与重复请求 & 会话恢复、请求编号复用、并发写入顺序和题目进度保护 \\
    语音与数字人 & 听写签名保护、术语整理、音频提交、数字人公开配置和透明画面尺寸 \\
    学习画像 & 问诊事件分类、分项掌握度、薄弱点和时间线筛选 \\
    \bottomrule
  \end{tabularx}
\end{table}

\section{资源生成测试}

资源生成测试覆盖讲义、Tutor 讲解、思维导图、练习题、阶段卷、PPT/slides 和音频。

首节讲义测试中，当前课程节点绑定 1 个知识节点和 2 条资料片段。生成结果包含 1 段 177 字的课程概述和 4 个主要知识点，分别涉及牙周炎的定义与疾病性质、牙周炎与牙龈病的区别、始动因素与全身促进因素，以及临床表现与疾病进展。

\begin{table}[H]
  \centering
  \caption{资源生成测试结果}
  \label{tab:resource-test}
  \begin{tabularx}{\linewidth}{p{0.27\linewidth}X}
    \toprule
    资源类型 & 实际结果 \\
    \midrule
    首节讲义 & 1 个知识节点、2 条资料片段、177 字概述、4 个主要知识点 \\
    Tutor 讲解 & 返回 \texttt{llm\_generated} 内容 \\
    思维导图 & 根据 5 个课程节点生成课程结构 \\
    练习题 & 根据课程内容和知识节点生成题目 \\
    阶段卷 & 完成题目、答案、批改和结果返回 \\
    PPT/slides & 单页资源文件生成成功，资源路径存在；两步流程能够合并为完整 deck \\
    音频 & 讲解音频生成链路返回有效资源状态 \\
    \bottomrule
  \end{tabularx}
\end{table}

\section{学习测评与学习状态测试}

学习测评测试检查题目生成、答案提交、自动批改、答案解析、错因识别、薄弱知识节点更新和补救建议生成。

在学习状态测试中，用户先提交 1 次错误答案，系统生成 1 条答案记录和至少 1 条掌握度记录；随后提交正确答案，答案记录增加到 2 条，掌握度记录保持存在且掌握度分数发生变化。课程节点在答题和答疑过程中保持不变。

\begin{table}[H]
  \centering
  \caption{学习测评与状态更新结果}
  \label{tab:assessment-test}
  \begin{tabularx}{\linewidth}{p{0.35\linewidth}X}
    \toprule
    检查项目 & 实际结果 \\
    \midrule
    题目生成 & 依据当前课程节点生成带选项和正确答案的题目 \\
    错误答案处理 & \texttt{is\_correct=False}，保存用户选项和正确选项 \\
    错误解析 & 返回选项错误分析、正确答案解释和误解概述 \\
    错误后进度 & 当前课程节点保持不变 \\
    错误后画像 & 薄弱点和错误模式发生更新 \\
    掌握度记录 & 错误答案后新增或更新掌握度记录 \\
    正确答案处理 & 答案记录总数增加到 2 条 \\
    掌握度变化 & 正确答案后掌握度分数发生变化 \\
    答案持久化 & 会话摘要中包含 \texttt{answers} 字段 \\
    补救建议 & 根据错误知识节点生成后续学习建议 \\
    \bottomrule
  \end{tabularx}
\end{table}

用户模型测试还验证了画像从初始状态到学习中状态的变化：初始画像为本科高年级、期末考试、偏好病例带入、标准节奏、暂无薄弱点；中途提问后增加“机制链条不完整”和“牙槽骨吸收”两个薄弱点；答错题后进一步增加“牙龈炎与牙周炎区别”；最终画像的偏好调整为“先病例后概念”，节奏调整为“慢一点”。

\section{知识检索与证据一致性测试}

真实 KAG 检索测试成功加载 34,805 条文本片段，同时加载稀疏倒排索引和 BGE-M3 稠密向量索引。以“牙周炎”为检索词时，系统返回 5 条相关结果，并展示前 3 条结果。结果来源包括牙周病学教材、口腔组织病理学教材和临床医学资料。

\begin{table}[H]
  \centering
  \caption{KAG 检索结果}
  \label{tab:kag-retrieval-test}
  \begin{tabularx}{\linewidth}{p{0.34\linewidth}X}
    \toprule
    检查项目 & 实际结果 \\
    \midrule
    文本片段加载 & 34,805 条 \\
    稀疏倒排索引 & 加载成功 \\
    BGE-M3 稠密索引 & 加载成功 \\
    “牙周炎”检索 & 返回 5 条结果 \\
    结果展示 & 展示前 3 条结果 \\
    资料来源 & 教材、医学手册和知识库资料 \\
    \bottomrule
  \end{tabularx}
\end{table}

课程节点的证据绑定和答疑过程均使用当前证据池中的合法标识。当前章节答疑使用 2 条证据；课程延伸答疑在定向检索后使用 4 条证据。非法知识节点或文本片段标识无法进入有效规划结果，模型返回的不同证据字段也能够被统一处理。

\section{响应时间测试}

响应时间测试使用本地后端真实学习接口完成，测试主题为“牙周炎”。

\begin{table}[H]
  \centering
  \caption{核心学习流程响应时间}
  \label{tab:latency-test}
  \begin{tabularx}{\linewidth}{p{0.40\linewidth}p{0.22\linewidth}X}
    \toprule
    测试阶段 & 实测耗时 & 实际结果 \\
    \midrule
    首节课程大纲就绪 & 145.088 秒 & 生成 5 个课程节点 \\
    首节讲义缓存就绪 & 随首节资源流程完成 & 缓存状态为 \texttt{ready} \\
    Tutor 首次内容就绪 & 195.103 秒 & 内容状态为 \texttt{llm\_generated} \\
    完整学习请求完成 & 195.103 秒 & 接口返回成功 \\
    \bottomrule
  \end{tabularx}
\end{table}

\section{稳定性测试}

稳定性测试覆盖计划草稿加载、计划项目依赖关系、学习会话隔离、学习状态保存与恢复、画像状态读取、检索数据路径、缓存状态和非法数据校验。

\begin{table}[H]
  \centering
  \caption{稳定性测试结果}
  \label{tab:stability-test}
  \begin{tabularx}{\linewidth}{p{0.42\linewidth}X}
    \toprule
    检查项目 & 实际结果 \\
    \midrule
    计划草稿加载 & 草稿标识能够解析，草稿项目能够映射为推荐卡片 \\
    计划项目依赖 & 多项目依赖按拓扑顺序处理，缺失依赖、负数依赖、自依赖和循环依赖被拦截 \\
    会话隔离 & 不同会话分别保存各自的计划草稿 ID 和后端会话 ID \\
    学习状态恢复 & 会话保存后能够恢复课程进度、答案记录和掌握度记录 \\
    画像状态读取 & 能够读取基础、目标、偏好、节奏和薄弱点字段 \\
    检索数据路径 & 能够定位本地嵌套数据目录下的知识库和索引文件 \\
    讲义缓存状态 & 首节讲义缓存状态返回 \texttt{ready} \\
    非法证据校验 & 不属于当前证据池的知识节点和文本片段不能进入有效规划结果 \\
    \bottomrule
  \end{tabularx}
\end{table}

\section{测试结果汇总}

\begin{table}[H]
  \centering
  \caption{测试结果汇总}
  \label{tab:test-summary}
  \begin{tabularx}{\linewidth}{p{0.25\linewidth}X}
    \toprule
    测试类别 & 主要结果 \\
    \midrule
    自动化测试 & 既有 961 项 Python 测试完成，数字问诊新增 49 项专项测试并全部通过 \\
    数字问诊 & 自然语言诊断、患者角色边界、双路证据、判官交接、语音术语、重复请求和画像写入通过专项验证 \\
    功能测试 & “牙周炎”主题生成 5 个课程节点，完成讲义、Tutor 和 QA 链路 \\
    学习画像 & 薄弱点从空集更新为 3 个具体薄弱方向，偏好和学习节奏发生状态更新 \\
    学习测评 & 错误答案后保存 1 条答案记录并更新掌握度；正确答案后答案记录增加到 2 条且掌握度分数变化 \\
    资源生成 & 首节讲义包含 177 字概述和 4 个知识点；PPT 单页文件和完整 deck 生成成功 \\
    KAG 检索 & 加载 34,805 条文本片段，检索“牙周炎”返回 5 条结果并展示前 3 条 \\
    证据一致性 & 当前章节答疑使用 2 条证据，延伸答疑使用 4 条证据，并生成 2 条定向检索计划 \\
    响应时间 & 大纲就绪 145.088 秒，Tutor 首次内容就绪 195.103 秒 \\
    稳定性 & 计划依赖、会话隔离、状态恢复、缓存状态和非法数据校验均完成对应行为验证 \\
    \bottomrule
  \end{tabularx}
\end{table}

综合测试结果表明，OralSage 已形成由知识检索、课程规划、资源生成、智能讲解、即时答疑、数字问诊、学习测评和状态更新组成的学习闭环。自动化测试、数字问诊专项测试、真实检索测试、智能体链路测试和核心学习接口测试完成了对应验证，系统主要功能能够按照预期运行。
